\section{Experimental Setup} \label{sec:app-setting}

\subsection{Implementation Details of Baseline Methods}

\subsubsection{Direct Score Predictor.}
We fine-tune Qwen2.5-Omni-7B on the ConvTTS training set, which consists of 4,579 audio samples. Given an input speech recording, the model is trained to directly output numerical ratings for all evaluation sub-metrics as well as the overall human-likeness score, without generating any intermediate chain-of-thought (CoT) reasoning. The supervised fine-tuning (SFT) prompt template used for this baseline is shown below.

We use a learning rate of $2\times10^{-5}$ with a cosine learning rate scheduler. The batch size is set to 32, the maximum sequence length is 4096, and the model is trained for up to 10 epochs. All experiments are conducted using the SWIFT framework~\cite{zhao2024swiftascalablelightweightinfrastructure}.
\begin{tcolorbox}[gray_box, title = {{Prompt Template: Direct Score Predictor}}]\footnotesize
You are an expert audio rater specializing in evaluating the quality of audio recordings.

You will be provided with an audio recording of a conversation between a TTS system and a user. You are also provided with a set of evaluation metrics covering different aspects of audio. 

Your task is to write a natural, clear, and concise assessment of each evaluation metric and assign a rating based on your listening experience. \\

\# Evaluation Metrics

Each evaluation metric is described below:

- expressive_intensity: The intensity of expressiveness in the audio.

- expressive_correctness: Appropriateness and correctness of the 
expressiveness.

- intonation: Quality of the intonation in the speech.

- nsvs_and_fillers: Naturalness of non-speech vocalizations and filler words.

- mispronunciation: whether mispronunciation is present.

- pacing: Naturalness of the pacing or speed of the speech.

- overall_score: Overall Human likeness rating

Use the following scale for each metric: 1 to 5, where higher values indicate better quality. Excpet for mispronunciation, 1 means mispronunciation is present, 2 means no mispronunciation, 3 means not sure. \\

\# Instruction and Response Format

Please follow these steps to complete the task. Your response must strictly follow this format:

[Final Ratings]

expressive_intensity: <your rating>

expressive_correctness: <your rating>

intonation: <your rating>

nsvs_and_fillers: <your rating>

mispronunciation: <your rating>

pacing: <your rating>

overall_score: <your rating>

\end{tcolorbox}

\subsubsection{Few-shot Speech Prompting.}
Following the setup of \citet{manakul2025audiojudge}, we randomly select five audio samples with human-likeness ratings spanning 1 to 5 as in-context learning (ICL) examples, ensuring coverage from low to high-quality speech. The corresponding audio clips are concatenated into a single audio sequence, with short segments of background silence inserted between clips. We prompt Gemini-2.5-Pro in voice mode using this stitched audio together with a few-shot prompt that includes the ground-truth human ratings for the five examples. The full prompt template is provided below.
\begin{tcolorbox}[gray_box, title = {{Prompt Template: Few-shot Speech Prompting}}]\footnotesize
You are an expert audio rater specializing in evaluating the quality of audio recordings.

You will be provided with an audio recording of five demonstration conversation segments between a TTS system and a user. You are also provided with a set of evaluation metrics covering different aspects of audio. 

Your task is to write a natural, clear, and concise assessment of each evaluation metric and assign a rating to the final conversation segment (the sixth segment).

Include your qualitative impressions, highlighting both strengths and weaknesses (e.g., drifts in pitch style across turns, or flatness) of the audio. 

Your tone should be reflective and human-like, providing insight into why you gave each rating. Each claim must be supported with evidence from the audio recording. \\

\# Evaluation Metrics

Each evaluation metric is described below:

- expressive_intensity: The intensity of expressiveness in the audio.

- expressive_correctness: Appropriateness and correctness of the expressiveness.

- intonation: Quality of the intonation in the speech.

- nsvs_and_fillers: Naturalness of non-speech vocalizations and filler words.

- mispronunciation: whether mispronunciation is present.

- pacing: Naturalness of the pacing or speed of the speech.

- overall_score: Overall Human likeness rating of the second speaker (speaker-2), take all the above metrics into account.

Use the following scale for each metric: 1-5. Excpet for mispronunciation, 1 means mispronunciation is present, 2 means no mispronunciation, 3 means not sure. \\

\# Response Format

Please follow these steps to complete the task. Your response must strictly follow this format:

[Reasoning regarding the five few-shot segments]

<Explain how the ground truth ratings in the prompt can be derived from the provided few-shot audio segments.>

[Summary Explaining the Ratings of the final test example]

<Detailed evidence log of the final test audio segment (the sixth segment), including detailed analysis of each sentence.>

<Reasoning for your rating of each evaluation metric>

<The rating of overall_score should take all factors into account.> \\

[Final Ratings]

expressive_intensity: <your rating>

expressive_correctness: <your rating>

intonation: <your rating>

nsvs_and_fillers: <your rating>

mispronunciation: <your rating>

pacing: <your rating>

overall_score: <your rating> \\

\# Important

- Please only judge the final conversation segment (the sixth segment), not the few-shot segments. The few-shot segments are only used as reference to help you understand the ground truth ratings. \\

\# Few-Shot Segment Examples

\#\# Conversation-1

\#\# Ratings:

\{Ground-truth Human Ratings\}

\#\# Conversation-2

\#\# Ratings:

\{Ground-truth Human Ratings\}

\#\# Conversation-3

\#\# Ratings:

\{Ground-truth Human Ratings\}

\#\# Conversation-4

\#\# Ratings:

\{Ground-truth Human Ratings\}

\#\# Conversation-5

\#\# Ratings:

\{Ground-truth Human Ratings\} \\

\# Now evaluate the final conversation segment in audio clip:
\end{tcolorbox}

\subsubsection{Few-shot Acoustic Feature Prompting.}
Instead of providing raw audio as context, we prompt a text-only LLM using explicitly extracted acoustic features. These features are obtained using the same vowel-level acoustic feature extraction pipeline described in Section~\ref{sec:gsrm} and illustrated in Figure~\ref{fig:teaser}. For each few-shot example, we concatenate the extracted acoustic features with the corresponding ground-truth human ratings, followed by the acoustic features of the test sample. The text LLM is then prompted to predict the ratings for the test sample. The prompt template used for acoustic feature prompting is shown below.
\begin{tcolorbox}[gray_box, title = {{Prompt Template: Few-shot Acoustic Feature Prompting}}]\footnotesize
You are an expert audio rater specializing in evaluating the quality of audio recordings by leveraging acoustic features.

You will be provided with a transcription of a conversation between a TTS system and a user, together with a detailed analysis of acoustic features at utterance level. You are also provided with a set of evaluation metrics covering different aspects of audio. 

Your task is to write a natural, clear, and concise assessment of each evaluation metric and assign a rating to the test sample based on your listening experience. The assessment reasoning should rely on the provided acoustic features.

Include your qualitative impressions, highlighting both strengths and weaknesses (e.g., drifts in pitch style across turns, or flatness) of the audio. Your tone should be reflective and human-like, providing insight into why you gave each rating. Each claim must be supported with evidence. \\

\# Evaluation Metrics

Each evaluation metric is described below:

- expressive_intensity: The intensity of expressiveness in the audio.

- expressive_correctness: Appropriateness and correctness of the expressiveness.

- intonation: Quality of the intonation in the speech.

- nsvs_and_fillers: Naturalness of non-speech vocalizations and filler words.

- mispronunciation: whether mispronunciation is present.

- pacing: Naturalness of the pacing or speed of the speech.

- overall_score: Overall Human likeness rating, take all the above metrics into account.

Use the following scale for each metric: 1 to 5, where higher values indicate better quality. Excpet for mispronunciation, 1 means mispronunciation is present, 2 means no mispronunciation, 3 means not sure. \\

\# Acoustic Features

<Detailed evidence log of the audio recording, including vowel-level analysis of each sentence.> \\

\# Instructions

When evaluating the audio, you must rely on the acoustic features—Pitch, Pitch Variation, Pitch Slope, Intensity, and Intensity Variation—as your primary evidence. Use them in the following expert-informed ways:

- Expressive Intensity: Identify the strength of expressiveness by looking for moderate to high Pitch Variation, clear Pitch Slopes, and sufficient Intensity. Very low variation in pitch or intensity indicates a flat, emotionless delivery. Extremely high variation may indicate unstable, exaggerated, or inconsistent expressiveness.

- Expressive Correctness: Evaluate whether the expressiveness is appropriate. Check if Pitch Variation and Intensity Variation are coherent and context-appropriate, not abrupt or mismatched. Excessively sharp slopes or sudden changes in loudness may suggest incorrect or unnatural expressiveness.

- Intonation Quality: Examine Pitch Slope to understand rising/falling contours, and Pitch Variation to assess natural melodic movement. Smooth slopes and moderate variation indicate good intonation. Flat pitch indicates monotone delivery; abrupt or irregular slopes indicate unstable intonation.

- NSVs and Fillers (Non-Speech Vocalizations): Look for sudden spikes in Intensity Variation or Pitch Variation that do not align with vowel structure. Irregular or abrupt changes may reflect throat noises, glottal catches, or filler-like disfluencies.

- Mispronunciation: Evaluate vowels for abnormally low intensity, very short or weak articulation (if duration provided), or unstable pitch contours. A vowel whose acoustic pattern deviates sharply from surrounding vowels may indicate mispronunciation.

- Pacing: Assess pacing consistency by checking whether the prosodic behavior across vowels is smooth and regular. Extreme or uneven Intensity Variation or Pitch. \\

\# Response Format

[Summary Explaining the Ratings]

<First check the above analysis for completeness and consistency; correct any errors; and add any missing issues or positives.>

<Reasoning for your rating of each evaluation metric. Refer to the details in the analysis.>

<The rating of overall_score should take all factors into account.> \\

[Final Ratings]

expressive_intensity: <your rating>

expressive_correctness: <your rating>

intonation: <your rating>

nsvs_and_fillers: <your rating>

mispronunciation: <your rating>

pacing: <your rating>

overall_score: <your rating> \\

\# Few-Shot Examples

\#\# Conversation-1

\#\# Acoustic Feature:

\{acoustic feature of sample-1\}

\#\# Ratings:

\{Ground-truth Human Ratings\}

\#\# Conversation-2

\#\# Acoustic Feature:

\{acoustic feature of sample-2\}

\#\# Ratings:

\{Ground-truth Human Ratings\}

\#\# Conversation-3

\#\# Acoustic Feature:

\{acoustic feature of sample-3\}

\#\# Ratings:

\{Ground-truth Human Ratings\}

\#\# Conversation-4

\#\# Acoustic Feature:

\{acoustic feature of sample-4\}

\#\# Ratings:

\{Ground-truth Human Ratings\}

\#\# Conversation-5

\#\# Acoustic Feature:

\{acoustic feature of sample-5\}

\#\# Ratings:

\{Ground-truth Human Ratings\} \\

\# Now evaluate the final conversation:

\#\# Acoustic Feature:

\{acoustic feature of test sample\}
\end{tcolorbox}


\subsubsection{MOS Predictors.}
We include existing utterance-level MOS predictors as baselines, specifically NISQA\footnote{The NISQA model code used in this research is licensed under the MIT License. The model weights are licensed under the Creative Commons Attribution-NonCommercial-ShareAlike 4.0 International (CC BY-NC-SA 4.0) License. The weights were used solely for non-commercial research purposes in accordance with the terms of this license. No distribution of the model weights or code containing the weights is included in this publication.}~\cite{mittag2021nisqa} and UTMOSv2~\cite{baba2024utmosv2}. Both models are originally designed for single-speaker, single-utterance speech naturalness assessment and output a scalar MOS score. Since these predictors do not support multi-turn conversational inputs, we apply them independently to each utterance within a conversation. For example, if a conversation consists of five utterances, the predictor produces five corresponding MOS scores. Conversation-level predictions are then obtained by aggregating the utterance-level scores using simple statistics, such as the mean or minimum.

\subsubsection{Regression-based Predictors.}
In contrast to MOS predictors, regression-based predictors are explicitly trained to map conversational speech representations to human-annotated MOS scores. These models follow a unified regression framework built on top of large pre-trained speech encoders, including AES- and WavLM-based backbones. Given an input conversation, the audio is segmented into fixed-length windows, and frame-level representations are extracted using a frozen transformer encoder. A weighted sum of hidden states across all encoder layers is computed, followed by temporal mean pooling to obtain a segment-level embedding. The resulting embeddings are then passed through a multi-layer perceptron with GeLU activations and dropout, trained with mean squared error loss to predict MOS scores. Final conversation-level predictions are obtained by averaging the regression outputs across all segments.


\subsection{Implementation Details of GSRM}

\subsubsection{Acoustic Feature Extraction.}
\label{app:acoustic_features}

Given a speech recording and its transcript, all audio is first resampled to 24 kHz. For multi-channel recordings, only the first channel is retained (system response). Extended silences are removed using a Silero-based voice activity detector with a merge threshold of 1000 ms to preserve natural intra-utterance timing.

Phoneme-level forced alignment is then performed between the transcript and waveform using a neural alignment model. From the resulting alignment, we retain only vowel segments. The vowel inventory is defined using an IPA-based set covering both monophthongs and diphthongs. For each aligned vowel, two temporal spans are defined. The \emph{full span}, computed from the earliest onset to the latest offset across all alignment states, is used to measure vowel duration. The \emph{core span} is a stable sub-region used for pitch and intensity estimation to avoid boundary transitions. When five or more alignment states are available, the middle two states are selected; when three to four states are present, the central state is used; otherwise, 20\% of the full span is trimmed from each edge.

From each vowel segment, we extract a compact set of low-level prosodic features: (i) pitch level, defined as the mean fundamental frequency (F0) over the core span; (ii) pitch variation, computed as the standard deviation of F0; (iii) pitch slope, obtained via a linear fit to F0 over time; (iv) intensity level, measured as mean RMS energy; (v) intensity variation, defined as the standard deviation of RMS energy; and (vi) duration, defined as the length of the full span. Pitch extraction follows a Praat-style autocorrelation method with speaker-adaptive floor and ceiling parameters, and all pitch and intensity statistics are computed only over voiced frames.

Continuous feature values are normalized using a two-stage procedure. First, features are z-normalized at the speaker level to control for individual vocal characteristics. Second, vowel-type normalization is applied to mitigate systematic differences across vowel classes. The normalized values are then discretized into ordinal categories. Specifically, pitch level is discretized using fixed thresholds at 85, 110, 160, 220, 280, and 390~Hz, corresponding to \emph{extremely low}, \emph{very low}, \emph{low}, \emph{mid}, \emph{high}, \emph{very high}, and \emph{extremely high}. Pitch slope is categorized separately for rising and falling contours using the 25th and 75th percentiles of positive and negative slopes, yielding \emph{low}, \emph{mid}, and \emph{high} slope categories. Pitch variation, intensity level, intensity variation, and duration are discretized via quantile-based binning at the 10th, 25th, 75th, and 90th percentiles, producing five ordinal categories: \emph{very low}, \emph{low}, \emph{mid}, \emph{high}, and \emph{very high}. Duration quantiles are computed only from vowels with fully voiced cores.

The final output is an utterance-level acoustic feature log that records vowel sequences and their discretized prosodic attributes in temporal order. We provide an example of the extracted vowel-level acoustic feature log below.

\begin{tcolorbox}[blue_box, title = {{Example: Vowel-level Acoustic Feature}}]\footnotesize
\#\# Transcript: ``so there could be a few reasons why he's doing this."

\#\#\# The vowel /\IPA{o}/ in ``so": [Pitch] mid, [Pitch Variation] low, [Pitch Slope] high falling, [Intensity] high, [Intensity variation] mid, [Duration] very high.

\#\#\# The vowel /\IPA{ɛ}/ in ``there": [Pitch] mid, [Pitch Variation] mid, [Pitch Slope] high falling, [Intensity] high, [Intensity variation] very low, [Duration] low.

\#\#\# The vowel /\IPA{ʊ}/ in ``could": [Pitch] high, [Pitch Variation] very high, [Pitch Slope] high falling, [Intensity] mid, [Intensity variation] mid, [Duration] low.

\#\#\# The vowel /\IPA{iː}/ in ``be": [Pitch] high, [Pitch Variation] mid, [Pitch Slope] high falling, [Intensity] mid, [Intensity variation] low, [Duration] mid.

\#\#\# The vowel /\IPA{ə}/ in ``a": [Pitch] high, [Pitch Variation] high, [Pitch Slope] high falling, [Intensity] high, [Intensity variation] mid, [Duration] low.

\#\#\# The vowel /\IPA{uː}/ in ``few": [Pitch] mid, [Pitch Variation] mid, [Pitch Slope] high falling, [Intensity] mid, [Intensity variation] mid, [Duration] mid.

\#\#\# The vowel /\IPA{iː}/ in ``reasons": [Pitch] mid, [Pitch Variation] mid, [Pitch Slope] high falling, [Intensity] mid, [Intensity variation] low, [Duration] mid.

\#\#\# The vowel /\IPA{a}/ in ``why": [Pitch] low, [Pitch Variation] mid, [Pitch Slope] high falling, [Intensity] mid, [Intensity variation] very low, [Duration] low.

\#\#\# The vowel /\IPA{iː}/ in ``he's": [Pitch] low, [Pitch Variation] mid, [Pitch Slope] high falling, [Intensity] low, [Intensity variation] mid, [Duration] mid.

\#\#\# The vowel /\IPA{uː}/ in ``doing": [Pitch] low, [Pitch Variation] mid, [Pitch Slope] high falling, [Intensity] mid, [Intensity variation] low, [Duration] mid.

\#\#\# The vowel /\IPA{ɪ}/ in ``this": [Pitch] low, [Pitch Variation] mid, [Pitch Slope] high falling, [Intensity] low, [Intensity variation] mid, [Duration] high.

\#\# Transcript: ``he might be seeking validation from his friends, or maybe he's not aware of how his behavior affects you."

\#\#\# The vowel /\IPA{iː}/ in ``he": [Pitch] high, [Pitch Variation] low, [Pitch Slope] low rising, [Intensity] mid, [Intensity variation] mid, [Duration] low.

\#\#\# The vowel /\IPA{a}/ in ``might": [Pitch] high, [Pitch Variation] mid, [Pitch Slope] mid rising, [Intensity] mid, [Intensity variation] very low.

\#\#\# The vowel /\IPA{iː}/ in ``be": [Pitch] high, [Pitch Variation] high, [Pitch Slope] high falling, [Intensity] mid, [Intensity variation] low, [Duration] low.

\#\#\# The vowel /\IPA{iː}/ in ``seeking": [Pitch] mid, [Pitch Variation] high, [Pitch Slope] high falling, [Intensity] mid, [Intensity variation] high, [Duration] mid.

\#\#\# The vowel /\IPA{æ}/ in ``validation": [Pitch] mid, [Pitch Variation] mid, [Pitch Slope] high falling, [Intensity] mid, [Intensity variation] low, [Duration] mid.

\#\#\# The vowel /\IPA{ɪ}/ in ``his": [Pitch] low, [Pitch Variation] mid, [Pitch Slope] high falling, [Intensity] low, [Intensity variation] mid, [Duration] low.

\#\#\# The vowel /\IPA{ɛ}/ in ``friends": [Pitch] low, [Pitch Variation] mid, [Pitch Slope] high falling, [Intensity] low, [Intensity variation] mid, [Duration] high.

\#\#\# The vowel /\IPA{ɔː}/ in ``or": [Pitch] low, [Pitch Variation] mid, [Pitch Slope] high falling, [Intensity] low, [Intensity variation] very high.

\#\#\# The vowel /\IPA{e}/ in ``maybe": [Pitch] low, [Pitch Variation] low, [Pitch Slope] high falling, [Intensity] low, [Intensity variation] very low, [Duration] mid.

\#\#\# The vowel /\IPA{iː}/ in ``maybe": [Pitch] low, [Pitch Variation] mid, [Pitch Slope] high falling, [Intensity] low, [Intensity variation] mid, [Duration] mid.

\#\#\# The vowel /\IPA{iː}/ in ``he's": [Pitch] low, [Pitch Variation] mid, [Pitch Slope] high falling, [Intensity] low, [Intensity variation] mid, [Duration] mid.

\#\#\# The vowel /\IPA{ɒ}/ in ``not": [Pitch] low, [Pitch Variation] mid, [Pitch Slope] high falling, [Intensity] low, [Intensity variation] very low, [Duration] mid.

\#\#\# The vowel /\IPA{ə}/ in ``aware": [Pitch] low, [Pitch Variation] low, [Pitch Slope] high falling, [Intensity] low, [Intensity variation] mid, [Duration] mid.

\#\#\# The vowel /\IPA{a}/ in ``how": [Pitch] low, [Pitch Variation] mid, [Pitch Slope] high falling, [Intensity] low, [Intensity variation] high, [Duration] mid.

\#\#\# The vowel /\IPA{ɪ}/ in ``his": [Pitch] low, [Pitch Variation] low, [Pitch Slope] mid rising, [Intensity] low, [Intensity variation] low, [Duration] mid.

\#\#\# The vowel /\IPA{ɪ}/ in ``behavior": [Pitch] low, [Pitch Variation] mid, [Pitch Slope] high falling, [Intensity] low, [Intensity variation] mid, [Duration] mid.

\#\#\# The vowel /\IPA{uː}/ in ``you": [Pitch] low, [Pitch Variation] mid, [Pitch Slope] high falling, [Intensity] very low, [Intensity variation] very high.

\#\# Transcript: ``he could be trying to assert dominance or control, or maybe he's got some insecurity issues."

\#\#\# The vowel /\IPA{iː}/ in ``he": [Pitch] high, [Pitch Variation] mid, [Pitch Slope] high falling, [Intensity] mid, [Intensity variation] high, [Duration] mid.

\#\#\# The vowel /\IPA{ʊ}/ in ``could": [Pitch] high, [Pitch Variation] mid, [Pitch Slope] high falling, [Intensity] mid, [Intensity variation] mid, [Duration] low.

\#\#\# The vowel /\IPA{iː}/ in ``be": [Pitch] high, [Pitch Variation] mid, [Pitch Slope] high falling, [Intensity] mid, [Intensity variation] mid, [Duration] mid.

\#\#\# The vowel /\IPA{a}/ in ``trying": [Pitch] mid, [Pitch Variation] mid, [Pitch Slope] high falling, [Intensity] mid, [Intensity variation] low, [Duration] low.

\#\#\# The vowel /\IPA{ə}/ in ``to": [Pitch] mid, [Pitch Variation] mid, [Pitch Slope] high falling, [Intensity] mid, [Intensity variation] mid, [Duration] low.

\#\#\# The vowel /\IPA{ə}/ in ``assert": [Pitch] mid, [Pitch Variation] mid, [Pitch Slope] high falling, [Intensity] mid, [Intensity variation] mid, [Duration] mid.

\#\#\# The vowel /\IPA{ɒ}/ in ``dominance": [Pitch] mid, [Pitch Variation] mid, [Pitch Slope] high falling, [Intensity] mid, [Intensity variation] low, [Duration] high.

\#\#\# The vowel /\IPA{ə}/ in ``control": [Pitch] low, [Pitch Variation] mid, [Pitch Slope] high falling, [Intensity] mid, [Intensity variation] low, [Duration] low.

\#\#\# The vowel /\IPA{ɔː}/ in ``or": [Pitch] mid, [Pitch Variation] mid, [Pitch Slope] low rising, [Intensity] very low, [Intensity variation] very high.

\#\#\# The vowel /\IPA{ɒ}/ in ``got": [Pitch] mid, [Pitch Variation] high, [Pitch Slope] high falling, [Intensity] mid, [Intensity variation] mid, [Duration] mid.

\#\#\# The vowel /\IPA{ʌ}/ in ``some": [Pitch] mid, [Pitch Variation] mid, [Pitch Slope] high falling, [Intensity] mid, [Intensity variation] low, [Duration] low.

\#\#\# The vowel /\IPA{ɪ}/ in ``insecurity": [Pitch] mid, [Pitch Variation] very low, [Pitch Slope] high falling, [Intensity] mid, [Intensity variation] very low, [Duration] low.

\#\#\# The vowel /\IPA{ɪ}/ in ``issues": [Pitch] low, [Pitch Variation] mid, [Pitch Slope] high falling, [Intensity] mid, [Intensity variation] low, [Duration] low.

\#\# Transcript: ``it's also possible he's not considering your feelings."

\#\#\# The vowel /\IPA{iː}/ in ``he's": [Pitch] high, [Pitch Variation] mid, [Pitch Slope] high falling, [Intensity] mid, [Intensity variation] low, [Duration] mid.

\#\#\# The vowel /\IPA{ɪ}/ in ``it's": [Pitch] mid, [Pitch Variation] very low, [Pitch Slope] low rising, [Intensity] very low, [Intensity variation] very high.

\#\#\# The vowel /\IPA{ɒ}/ in ``possible": [Pitch] mid, [Pitch Variation] high, [Pitch Slope] high falling, [Intensity] mid, [Intensity variation] high, [Duration] high.

\#\#\# The vowel /\IPA{iː}/ in ``he's": [Pitch] mid, [Pitch Variation] mid, [Pitch Slope] high falling, [Intensity] mid, [Intensity variation] mid, [Duration] high.

\#\#\# The vowel /\IPA{ɒ}/ in ``not": [Pitch] mid, [Pitch Variation] low, [Pitch Slope] high falling, [Intensity] mid, [Intensity variation] high, [Duration] high.

\#\#\# The vowel /\IPA{ə}/ in ``considering": [Pitch] mid, [Pitch Variation] mid, [Pitch Slope] high falling, [Intensity] mid, [Intensity variation] mid.

\#\#\# The vowel /\IPA{iː}/ in ``feelings": [Pitch] mid, [Pitch Variation] low, [Pitch Slope] high falling, [Intensity] mid, [Intensity variation] low, [Duration] high.

\#\# Transcript: ``have you talked to him about how you feel when this happens?"

\#\#\# The vowel /\IPA{æ}/ in ``have": [Pitch] mid, [Pitch Variation] low, [Pitch Slope] high falling, [Intensity] mid, [Intensity variation] mid, [Duration] mid.

\#\#\# The vowel /\IPA{ɔː}/ in ``talked": [Pitch] high, [Pitch Variation] mid, [Pitch Slope] high falling, [Intensity] mid, [Intensity variation] mid, [Duration] mid.

\#\#\# The vowel /\IPA{ɪ}/ in ``him": [Pitch] high, [Pitch Variation] mid, [Pitch Slope] high falling, [Intensity] mid, [Intensity variation] very low.

\#\#\# The vowel /\IPA{ə}/ in ``about": [Pitch] mid, [Pitch Variation] high, [Pitch Slope] high falling, [Intensity] mid, [Intensity variation] mid, [Duration] low.

\#\#\# The vowel /\IPA{a}/ in ``how": [Pitch] mid, [Pitch Variation] mid, [Pitch Slope] high falling, [Intensity] mid, [Intensity variation] mid, [Duration] mid.

\#\#\# The vowel /\IPA{uː}/ in ``you": [Pitch] mid, [Pitch Variation] mid, [Pitch Slope] high falling, [Intensity] low, [Intensity variation] mid, [Duration] mid.

\#\#\# The vowel /\IPA{iː}/ in ``feel": [Pitch] mid, [Pitch Variation] mid, [Pitch Slope] high falling, [Intensity] mid, [Intensity variation] low, [Duration] mid.

\#\#\# The vowel /\IPA{ɛ}/ in ``when": [Pitch] mid, [Pitch Variation] mid, [Pitch Slope] high falling, [Intensity] mid, [Intensity variation] very low, [Duration] low.

\#\#\# The vowel /\IPA{ɪ}/ in ``this": [Pitch] low, [Pitch Variation] mid, [Pitch Slope] high falling, [Intensity] low, [Intensity variation] mid, [Duration] low.

\#\#\# The vowel /\IPA{æ}/ in ``happens": [Pitch] mid, [Pitch Variation] mid, [Pitch Slope] high falling, [Intensity] low, [Intensity variation] mid, [Duration] mid.

\end{tcolorbox}


\subsubsection{Reasoning Synthesis.}
After extracting the acoustic feature log for an audio sample, we leverage a teacher model, GPT-4o, to synthesize the utterance-level evidence log. Specifically, for each utterance (sentence) and its associated vowel-level acoustic features, we prompt GPT-4o to generate a structured analysis capturing the inferred context, perceptual strengths, and potential issues of the utterance. We then concatenate the analyses of all utterances to form the complete evidence log for the audio sample. The prompt template used for evidence log synthesis is provided below.
\begin{tcolorbox}[gray_box, title = {{Prompt Template: Synthesizing Utterence-level Evidence Log}}]\footnotesize
You are an expert audio rater specializing in evaluating the quality of audio recordings. Your are given,

- A transcript of a spoken sentence.

- Acoustic features describing each vowel in the sentence.

- Evaluation metrics you must judge.

Your goal is to produce an analysis:

[Inferred context]: 

question/empathetic/apology/celebration/instruction/neutral statement. If no clear context exists, say: “No clear context inferred.”

[NSVs/fillers]: 

List any non-speech vocalizations (NSVs) or filler words (laugh, yeah, Aww, uh, hmm, etc.) that are present in the audio. If none exist, say: “No NSVs/fillers detected.”

[Positive]: 

Highlight the most positive aspects of the audio: expressive intensity, good emphasis, varied intonation, natural fillers/laughs, pleasant timbre. If no clear positive aspects exist, say: “No positive aspects detected.”

[Potential Issue]: 

Point out any potential issues: sharp pitch changes, intonation problems, pacing issues, speaker style changes, expressive incorrectness, mispronunciation, flat intonation, unnatural pauses, or any other issues. If no clear issue exists, say: “No notable issues detected.” \\

\# Evaluation Metrics

Each evaluation metric is described below:

- expressive_intensity: The intensity of expressiveness in the audio.

- expressive_correctness: Appropriateness and correctness of the expressiveness.

- intonation: Quality of the intonation in the speech.

- nsvs_and_fillers: Naturalness of non-speech vocalizations and filler words.

- mispronunciation: whether mispronunciation is present.

- pacing: Naturalness of the pacing or speed of the speech. \\

\# The system responses and acoustic features:

\{Acoustic features\} \\

\# Instructions

When evaluating the audio, you must rely on the provided acoustic features—Pitch, Pitch Variation, Pitch Slope, Intensity, and Intensity Variation—as your primary evidence. 

Use them in the following expert-informed ways:

- Expressive Intensity: Identify the strength of expressiveness by looking for moderate to high Pitch Variation, clear Pitch Slopes, and sufficient Intensity. Very low variation in pitch or intensity indicates a flat, emotionless delivery. Extremely high variation may indicate unstable, exaggerated, or inconsistent expressiveness.

- Expressive Correctness: Evaluate whether the expressiveness is appropriate. Check if Pitch Variation and Intensity Variation are coherent and context-appropriate, not abrupt or mismatched. Excessively sharp slopes or sudden changes in loudness may suggest incorrect or unnatural expressiveness.

- Intonation Quality: Examine Pitch Slope to understand rising/falling contours, and Pitch Variation to assess natural melodic movement. Smooth slopes and moderate variation indicate good intonation. Flat pitch indicates monotone delivery; abrupt or irregular slopes indicate unstable intonation.

- NSVs and Fillers (Non-Speech Vocalizations): Look for sudden spikes in Intensity Variation or Pitch Variation that do not align with vowel structure. Irregular or abrupt changes may reflect throat noises, glottal catches, or filler-like disfluencies.

- Mispronunciation: Evaluate vowels for abnormally low intensity, very short or weak articulation (if duration provided), or unstable pitch contours. A vowel whose acoustic pattern deviates sharply from surrounding vowels may indicate mispronunciation.

- Pacing: Assess pacing consistency by checking whether the prosodic behavior across vowels is smooth and regular. Extreme or uneven Intensity Variation or Pitch Variation may reflect rushed or unnatural pacing. If duration is given, long durations imply slow pacing; short durations imply fast pacing. \\

\# Response Format

Your response must strictly follow this format:

[Inferred context]

<one concise sentence describing the inferred context, referencing the transcripts as evidence.> OR ``No clear context inferred."

[NSVs/fillers]

<one concise sentence listing any NSVs/fillers, referencing the transcripts and features as evidence.> OR ``No NSVs/fillers detected."

[Positive]

<one concise sentence highlighting positive aspects of the evaluation metrics, referencing the features as evidence.> OR ``No positive aspects detected."

[Potential Issue]

<one concise sentence describing weaknesses, referencing the features as evidence.> OR ``No notable issues detected." \\

\# Important

- Always refer specifically to the words and vowels in the transcript as evidence. Do Not say something too generic without referencing the features.

- Although always use features and transcripts as evidence, Do Not explicitly mention something like ``Based on the features and transcripts, I judge the audio as follows:" in your response.
\end{tcolorbox}

Once the evidence log is constructed, we prompt GPT-4o again to synthesize a global judgment chain-of-thought (CoT). We adopt a \emph{role-playing} prompt that grants the model access to both the evidence log and the ground-truth human ratings, and instructs it to explain the reasoning behind the assigned scores. This design ensures that the synthesized reasoning is faithful to the oracle ratings while remaining coherent. The prompt template for global judgment synthesis is shown below.
\begin{tcolorbox}[gray_box, title = {{Prompt Template: Synthesizing Global Judgment CoT}}]\footnotesize
You are an audio assessment assistant striving to become an expert audio rater specializing in evaluating the quality of audio recordings. You are now taking an exam to evaluate your capabilities. 

You will be provided with a transcription of a conversation between a TTS system and a user, together with a detailed analysis of acoustic features at sentence level (user turns are ommitted). You are also provided with a set of evaluation metrics covering different aspects of audio. 

Your task is to write a natural, clear, and concise assessment of each evaluation metric and assign a rating based on your listening experience. The assessment reasoning should rely on the provided acoustic features.

Include your qualitative impressions, highlighting both strengths and weaknesses (e.g., drifts in pitch style across turns, or flatness) of the audio. Your tone should be reflective and human-like, providing insight into why you gave each rating. Each claim must be supported with evidence.

To evaluate your correctness, the oracle ratings will also be provided. You must ensure that your ratings MATCH the oracle ratings EXACTLY. 
However, your reasoning process must appear fully self-derived and must NOT reference, suggest awareness of, or appear to be influenced by the oracle ratings. \\

\# Evaluation Metrics

Each evaluation metric is described below:

- expressive_intensity: The intensity of expressiveness in the audio.

- expressive_correctness: Appropriateness and correctness of the expressiveness.

- intonation: Quality of the intonation in the speech.

- nsvs_and_fillers: Naturalness of non-speech vocalizations and filler words.

- mispronunciation: whether mispronunciation is present.

- pacing: Naturalness of the pacing or speed of the speech.

- overall_score: Overall Human likeness rating \\

\# Oracle Ratings (For Evaluation Only):

Use the following scale for each metric: 1-5. Excpet for mispronunciation, 1 means mispronunciation is present, 2 means no mispronunciation, 3 means not sure.

expressive_intensity: {expressive_intensity}

expressive_correctness: {expressive_correctness}

intonation: {intonation}

nsvs_and_fillers: {nsvs_and_fillers}

mispronunciation: {mispronunciation}

pacing: {pacing}

overall_score: {human_likeness_speaker_2} \\

\# The system responses and detailed analysis of acoustic features:

\{Evidence Log\} \\

\# Response Format 

Your response must strictly follow this format:

[Summary Explaining the Ratings]

<First check the above analysis for completeness and consistency; correct any errors; and add any missing issues or positives.>

<Reasoning for your rating of each evaluation metric. Refer to the details in the analysis.>

<The rating of overall_score should take all factors into account.> \\

[Final Ratings]

expressive_intensity: <your rating>

expressive_correctness: <your rating>

intonation: <your rating>

nsvs_and_fillers: <your rating>

mispronunciation: <your rating>

pacing: <your rating>

overall_score: <your rating>
\end{tcolorbox}

\subsubsection{Training and Inference.}
We fine-tune GSRM on the ConvTTS training set using Qwen2.5-Omni-7B as the backbone model. Given an input speech recording, GSRM is trained to first generate the utterance-level evidence log, followed by the global judgment CoT and the final rating scores. The supervised fine-tuning (SFT) prompt template used for training GSRM is provided below.
\begin{tcolorbox}[gray_box, title = {{Prompt Template: GSRM Training \& Inference }}]\footnotesize
You are an expert audio rater specializing in evaluating the quality of audio recordings.

You will be provided with an audio recording of a conversation between a TTS system and a user. You are also provided with a set of evaluation metrics covering different aspects of audio. 

Your task is to write a natural, clear, and concise assessment of each evaluation metric and assign a rating based on your listening experience.

Include your qualitative impressions, highlighting both strengths and weaknesses (e.g., drifts in pitch style across turns, or flatness) of the audio. 

Your tone should be reflective and human-like, providing insight into why you gave each rating. Each claim must be supported with evidence from the audio recording. \\

\# Evaluation Metrics

Each evaluation metric is described below:

- expressive_intensity: The intensity of expressiveness in the audio.

- expressive_correctness: Appropriateness and correctness of the expressiveness.

- intonation: Quality of the intonation in the speech.

- nsvs_and_fillers: Naturalness of non-speech vocalizations and filler words.

- mispronunciation: whether mispronunciation is present.

- pacing: Naturalness of the pacing or speed of the speech.

- overall_score: Overall Human likeness rating \\

\# Instruction and Response Format

Please follow these steps to complete the task. Your response must strictly follow this format: \\

[Evidence log]

<Detailed evidence log of the audio recording, including vowel-level analysis of each sentence.> \\

[Summary Explaining the Ratings]

<Reasoning for your rating of each evaluation metric>

<The rating of overall_score should take all factors into account.> \\

[Final Ratings]

expressive_intensity: <your rating>

expressive_correctness: <your rating>

intonation: <your rating>

nsvs_and_fillers: <your rating>

mispronunciation: <your rating>

pacing: <your rating>

overall_score: <your rating>
\end{tcolorbox}

We use a learning rate of $2\times10^{-5}$ with a cosine learning rate scheduler. The batch size is set to 32, the maximum sequence length is 4096, and the model is trained for up to 10 epochs. During inference, we select the model checkpoint with the lowest validation loss on a held-out validation set. For each audio sample, we perform up to 16 random inference passes with a sampling temperature of 1.0 and top-$p$ set to 0.6. The final rating is obtained by averaging the scores across these samples. All experiments are conducted using the SWIFT framework~\cite{zhao2024swiftascalablelightweightinfrastructure}.
